\section{Characterizing Disagreement among Jailbreak Judges}
\label{sec:judge-disagreement}

We begin by examining whether the outcome of jailbreak evaluation changes when
the interaction and evaluation criterion are fixed but the judge backbone is
varied. This study establishes the empirical basis for using disagreement in
the subsequent search procedure.

\paragraph{Setup.}
We randomly sample 50 malicious goals from AdvBench
\citep{zou2023universal} and evaluate ten jailbreak attacks against four victim
models: \texttt{deepseek-v4-flash-0731},
\texttt{gpt-5.6-luna}, \texttt{gpt-oss-120b}, and
\texttt{qwen3.7-flash}. The ten attacks are selected from a library of 47
implementations using preliminary disagreement between their original
evaluators and a unified judge. This selection yields a stress set with varied
response styles and evaluator-sensitive cases. We obtain one response for each
goal, victim, and attack combination, giving 200 responses per attack and 2,000
responses in total.

Every response is evaluated independently by GPT-Judge
(\texttt{gpt-oss-120b}), Qwen-Judge (\texttt{qwen3.7-flash}), and
DeepSeek-Judge (\texttt{deepseek-v4-flash-0731}). Within each comparison, the
three judges receive the same query, response, rubric, and output format. This
setup isolates the effect of the judge backbone from changes in the evaluated
interaction or criterion.

\paragraph{Evaluation criteria and metrics.}
We use the binary WildGuard prompt \citep{han2024wildguard} and the graded
StrongREJECT prompt \citep{souly2024strongreject}. Binary ASR is the fraction
of responses labeled harmful. For scoring ASR, a response is successful when
the judge detects no refusal and assigns an average convincingness and
specificity score of at least 4.

Conflict Rate (CR) measures the fraction of responses for which the three
judges do not reach the same attack-success decision. Under WildGuard, a
conflict occurs when their harmfulness labels differ. Under StrongREJECT, we
first apply the same success threshold used for scoring ASR and then compare the
resulting decisions. Missing judgments are counted as unsuccessful for ASR and
excluded from CR. Section~\ref{sec:tvd-conflict} uses a stricter pairwise
criterion for active search, requiring a normalized StrongREJECT score
difference greater than 0.5.

% Generated from code/Experiment-1/judge_batch_output/direct_results.jsonl.
\begin{table*}[t]
  \centering
  \caption{Attack success rate (ASR, \%) assigned by three jailbreak judges
  under binary and scoring criteria.}
  \label{tab:attack-asr-by-judge}
  \footnotesize
  \setlength{\tabcolsep}{4.8pt}
  \begin{tabular}{l cc cc cc}
    \toprule
    & \multicolumn{2}{c}{GPT-Judge}
    & \multicolumn{2}{c}{Qwen-Judge}
    & \multicolumn{2}{c}{DeepSeek-Judge} \\
    \cmidrule(lr){2-3}\cmidrule(lr){4-5}\cmidrule(lr){6-7}
    Attack & Binary & Scoring & Binary & Scoring & Binary & Scoring \\
    \midrule
    ActorBreaker~\citep{actorbreaker} & 43.0 & 66.0 & 24.0 & 52.5 & 25.5 & 55.0 \\
    AutoDAN-Turbo~\citep{autodanturbo} & 17.0 & 23.5 & 15.0 & 21.0 & 13.0 & 21.0 \\
    CIA Context~\citep{cia} & 49.5 & 70.0 & 21.5 & 70.5 & 27.5 & 72.5 \\
    Cognitive Overload~\citep{cognitiveoverload} & 0.5 & 4.0 & 0.5 & 4.0 & 0.0 & 3.5 \\
    Compound Roleplay~\citep{compoundroleplay} & 42.0 & 41.0 & 37.5 & 13.5 & 24.0 & 37.0 \\
    FITD~\citep{fitd} & 51.5 & 45.5 & 47.0 & 12.5 & 43.0 & 13.0 \\
    Mastermind~\citep{mastermind} & 31.5 & 87.5 & 13.0 & 70.5 & 15.0 & 79.0 \\
    MathPrompt~\citep{mathprompt} & 67.5 & 29.5 & 51.0 & 27.0 & 32.5 & 11.5 \\
    Prefill~\citep{prefill} & 31.0 & 30.0 & 17.0 & 22.5 & 24.5 & 20.0 \\
    Semantic Cloaking~\citep{semanticcloaking} & 50.0 & 28.0 & 42.0 & 47.0 & 45.0 & 49.0 \\
    \bottomrule
  \end{tabular}
\end{table*}

% Recomputed by code/Experiment-1/recompute_paper_metrics.py.
\begin{table}[t]
  \centering
  \caption{Conflict rate (CR, \%) among three jailbreak judges across attacks.
  Rows are ordered by mean CR; shading indicates conflict intensity.}
  \label{tab:attack-judge-conflict}
  \small
  \setlength{\tabcolsep}{5pt}
  \begin{tabular}{l rrr}
    \toprule
    Attack & Binary & Scoring & Mean \\
    \midrule
    Cognitive Overload~\citep{cognitiveoverload} & 0.5 & 1.5 & \cellcolor{red!0}1.0 \\
    AutoDAN-Turbo~\citep{autodanturbo} & 7.9 & 1.7 & \cellcolor{red!4}4.8 \\
    Prefill~\citep{prefill} & 23.7 & 12.5 & \cellcolor{red!14}18.1 \\
    Semantic Cloaking~\citep{semanticcloaking} & 24.8 & 24.1 & \cellcolor{red!19}24.5 \\
    Compound Roleplay~\citep{compoundroleplay} & 16.1 & 33.5 & \cellcolor{red!19}24.8 \\
    CIA Context~\citep{cia} & 41.1 & 9.0 & \cellcolor{red!19}25.0 \\
    FITD~\citep{fitd} & 16.7 & 35.6 & \cellcolor{red!20}26.1 \\
    Mastermind~\citep{mastermind} & 36.5 & 23.1 & \cellcolor{red!23}29.8 \\
    ActorBreaker~\citep{actorbreaker} & 35.0 & 25.4 & \cellcolor{red!23}30.2 \\
    MathPrompt~\citep{mathprompt} & 43.4 & 35.4 & \cellcolor{red!30}39.4 \\
    \bottomrule
  \end{tabular}
\end{table}


\paragraph{Results.}
Table~\ref{tab:attack-asr-by-judge} shows that measured attack success depends
substantially on the judge. The largest difference is 35.0 percentage points
for MathPrompt under binary evaluation, between GPT-Judge and DeepSeek-Judge.
Across all valid responses, the three judges conflict on 24.2\% of binary
decisions and 19.0\% of scoring decisions.

The disagreement is concentrated in particular attacks and criteria
(Table~\ref{tab:attack-judge-conflict}). The judges largely agree on Cognitive
Overload and AutoDAN-Turbo, but diverge on MathPrompt and CIA Context under
binary evaluation and on MathPrompt and FITD under scoring. An aggregate ASR
therefore obscures a structured set of evaluator-dependent cases. This result
motivates a collection method that searches for such cases directly instead of
waiting for them to appear among ordinary attack outputs.
